\documentclass[11pt,a4paper,twoside,openright]{book}
\usepackage[ngerman]{babel}
\newcommand{\HandbookPdfTitle}{SASD Math Toolkit - C++ Handbuch Version 0.4}
\newcommand{\HandbookFooter}{C++-Handbuch Version 0.4 - Deutsch}
\usepackage{sasd-cpp-handbook}

\begin{document}
\frontmatter

\begin{titlepage}
\pagecolor{sasdInk}\color{white}
\vspace*{2cm}
{\sffamily\bfseries\Large SASD-GmbH\par}
\vspace{3cm}
{\bfseries\fontsize{34}{40}\selectfont SASD Math Toolkit\par}
\vspace{0.5cm}
{\color{sasdTeal}\bfseries\fontsize{24}{29}\selectfont Modern C++\par}
\vspace{1cm}
{\Large Benutzer- und Entwicklerhandbuch\par}
{\large Version 0.4 / CPP-M3\par}
\vfill
{\large Portable Numerik, explizite Diagnostik und kontrollierte Mehrkernnutzung.\par}
\vspace{0.8cm}
{\repo\par}
\end{titlepage}
\nopagecolor\color{sasdInk}

\chapter*{Dokumentstatus}
\addcontentsline{toc}{chapter}{Dokumentstatus}

Dieses eigenständige C++-Handbuch beschreibt die CPP-M3-Version
\texttt{0.4.0}. Die Kapitel werden aus den gepflegten Markdown-Quellen unter
\texttt{docs/user-guide/cpp/de/} erzeugt. Es ist weder eine Kopie noch ein
Anhang des C\#/.NET-Handbuchs.

\begin{tabularx}{\textwidth}{@{}>{\bfseries}p{0.32\textwidth}X@{}}
Ausgabe & Modern C++20\\
Handbuchversion & 0.4\\
Bibliotheksversion & \texttt{0.4.0}\\
Root-Namespace & \texttt{sasd::numerics}\\
Buildsystem & CMake ab 3.25\\
Lizenz & MIT\\
Autor / Projekt & Robin Goerlach; SASD-GmbH\\
\end{tabularx}

\tableofcontents
\clearpage
\mainmatter

\input{build/de/README.tex}
\input{build/de/erste-schritte.tex}
\input{build/de/parallele-ausfuehrung.tex}
\input{build/de/deskriptive-statistik.tex}
\input{build/de/polynome-nullstellensuche.tex}
\input{build/de/interpolation-splines.tex}

\backmatter
\chapter*{Lizenz und Herkunft}
\addcontentsline{toc}{chapter}{Lizenz und Herkunft}
SASD Math Toolkit steht unter der MIT-Lizenz. Die C++-Implementierung und dieses
Handbuch entstehen unabhängig aus mathematischen Definitionen und den
sprachneutralen SASD-Verträgen; historischer Borland-Quellcode und historische
Handbuchtexte werden nicht übernommen.

\end{document}
